\section{Caractérisation physique de la latence de transport}

Afin de valider les performances du pont de communication (SPI $\rightarrow$ STM32 $\rightarrow$ Wi-Fi HaLow), une campagne de mesures physiques a été menée. L'objectif est d'isoler la latence de transmission pure, indépendamment des délais d'encodage vidéo et de capture du capteur.

\subsection{Méthodologie : Mesure à référence temporelle partagée}
Le principal défi de la mesure de latence réseau réside dans la synchronisation des horloges entre l'émetteur et le récepteur. Pour s'en affranchir, une architecture à "déclencheur commun" a été mise en place :
\begin{itemize}
    \item Deux Raspberry Pi 4B sont utilisées (une en TX, une en RX).
    \item Les masses (GND) des deux cartes sont reliées.
    \item Un bouton poussoir est connecté simultanément sur la broche \texttt{BCM 17} (Pin 11) des deux cartes, avec une résistance de rappel (\textit{pull-down}) de 10 k$\Omega$.
\end{itemize}

Lorsqu'on presse le bouton, les deux programmes détectent le front montant au même instant physique. Le TX envoie alors un paquet de test, et le RX arrête son chronomètre à la réception du paquet UDP.

\subsection{Implémentation des outils de test}
Des programmes spécifiques en C ont été développés pour exploiter cette architecture. Ils utilisent la bibliothèque \texttt{mmap} sur \texttt{/dev/gpiomem} pour garantir une lecture des entrées/sorties avec un minimum de gigue logicielle.

\begin{description}
    \item[Programme TX (\texttt{latency\_tx.c}) :] Attend l'appui sur le bouton, vérifie l'état du signal de \textit{handshake} du STM32 (Pin 18), puis injecte immédiatement un paquet de test de 1400 octets via le bus SPI.
    \item[Programme RX (\texttt{latency\_rx.c}) :] Déclenche un chronomètre haute résolution (\texttt{CLOCK\_MONOTONIC\_RAW}) dès la détection du bouton et le fige à l'arrivée du paquet sur le port UDP 1337 via un thread d'écoute dédié.
\end{description}

\subsection{Résultats des mesures de transport}
Deux types de tests ont été réalisés pour simuler un comportement unitaire et un comportement en flux continu :
\begin{enumerate}
    \item \textbf{Tests unitaires :} 10 envois espacés. La latence moyenne mesurée est de \textbf{11.06 ms}.
    \item \textbf{Test en rafale (\textit{Burst}) :} Envoi instantané de 10 paquets. La latence totale pour la rafale est de 124 ms, soit une moyenne en flux de \textbf{12.47 ms} par paquet.
\end{enumerate}

Ces résultats sont excellents : ils prouvent que le transport radio Sub-1 GHz ne représente qu'une fraction mineure (environ 6\%) de la latence totale observée visuellement.

\subsection{Isolation de la latence de capture}
Un test complémentaire de "boucle locale" (\textit{loopback}) a été effectué sur la Raspberry Pi 4B : capturer le flux caméra et l'afficher directement sur le même écran via GStreamer. 
La latence mesurée pour cette opération est de \textbf{5 à 7 ms}. 

\subsection{Conclusion du benchmarking}
Le cumul de la capture locale (~7 ms) et du transport radio (~12 ms) ne totalise que 19 ms. Or, la latence "Glass-to-Glass" finale constatée est supérieure à 150 ms. 
Cette divergence permet de conclure que les \textbf{130 ms restants} sont introduits par la pile 
logicielle de composition d'affichage de l'OS (X11/Wayland) ou par les algorithmes de mise en mémoire 
tampon des décodeurs vidéo. L'effort d'optimisation doit donc se concentrer sur le mode d'affichage 
"Bare-Metal" (\texttt{kmssink}).